\section{Integration patterns: how surrogates plug into ESM optimizers}
\label{sec:integration}

Section~\ref{sec:taxonomy} classifies surrogate families by
statistical structure. This section classifies how those families
attach to the host optimizer. The distinction matters because the
same neural network can either replace a market-clearing call inside
a mixed-integer programme (Pattern P1) or predict the programme's
solution in one shot (Pattern P3). Table~\ref{tab:T2-task-role} maps
patterns to task classes; Table~\ref{tab:T5-integration-patterns}
lists solver layer, safety practices and typical failure modes.
The subsections below follow P1--P5 and focus on multi-energy and
coupled-carrier studies in the screened pool; single-carrier
electricity examples appear only where no multi-carrier counterpart
is available in the pool.
Section~\ref{sec:doe} already uses many of the same papers to
illustrate training-data design; here we rotate the citation pool
toward studies whose primary contribution is the optimizer interface,
and cross-reference Section~\ref{sec:doe} when training and
integration are jointly documented.
The pattern subsections below describe every reference in
Tables~\ref{tab:T2-task-role} and~\ref{tab:T5-integration-patterns};
where a study is already used to illustrate training-data design in
Section~\ref{sec:doe}, the integration role is stated here and the
sampling workflow is not repeated.

% BEGIN inlined table table_T2_task_role_matrix
\begin{table*}[!t]
  \centering
  \scriptsize
  \caption{Energy-system task class (rows) versus surrogate integration
  role (columns); roles match patterns P1--P5 in
  Section~\ref{sec:integration}. Cells list the typical surrogate
  families and a small number of representative references; ``$\star$''
  marks the dominant pattern per row, ``--'' marks combinations that
  are rare in the screened pool.}
  \label{tab:T2-task-role}
  \renewcommand{\arraystretch}{1.25}
  \begin{tabularx}{\textwidth}{@{}L{0.16\textwidth} L{0.18\textwidth} L{0.16\textwidth} L{0.16\textwidth} L{0.14\textwidth} L{0.14\textwidth}@{}}
    \toprule
    Task class & P1 Replace subroutine & P2 Accelerate search & P3 Warm start / proxy & P4 Decomposition & P5 Uncertainty \\
    \midrule
    Economic dispatch / UC
      & $\star$ NN / L2O \cite{Wu2022231,Chen20244723,Pang2023}
      & GP-BO / kriging-NSGA \cite{Nikolaidis2021,Lin2023}
      & End-to-end NN \cite{Chen2022,Wu2022231}
      & Lagrangian / Benders \cite{Wu2024391,Zhang2025__4}
      & PCE / GP / quantile \cite{Hu2020,Hu2021671,Chen2023657} \\
    Optimal power flow
      & $\star$ NN / GNN \cite{Chen20244723,Chen2022,Liu20223726}
      & Kriging-EA \cite{Deng2020831,Lin2023}
      & L2O \cite{Li2023__2}
      & --
      & PCE / GP / DQ \cite{Muhlpfordt20194806,Chen2023657,Koirala20242284} \\
    Capacity expansion
      & NN / GBDT \cite{Jahangiri20244849,Pérez-Uresti2023}
      & GP-BO / NSGA \cite{Han2018622,Amaral20251156}
      & --
      & Benders + GP \cite{Zhang2025__4}
      & $\star$ Adaptive metamodels \cite{Fu20202904,Kim2021} \\
    District-heating dispatch
      & $\star$ NN metamodel \cite{Sánchez-Zabala2024,Du2025}
      & RSM / RFC \cite{Schulte2016104,Petrova2025}
      & ANN setpoints \cite{Reynolds2017704}
      & --
      & GP / NMPC \cite{Simonson202036} \\
    Multi-energy planning
      & NN / GP \cite{Zheng20231907,Yin2024,Roth2019214,Gao2025}
      & $\star$ GP-BO / NSGA \cite{Han2018622,Schulte2020,Aghaei Pour2022303}
      & Market-NN proxy \cite{Jalving2023,Wang2020202933}
      & --
      & PCE / GP \cite{Dong2023,Dong2022,Wang2024} \\
    Microgrid / energy hub
      & RBF / NN \cite{Wang2017,Hussien20211883,Li2023__3}
      & $\star$ SAEA / kriging-EA \cite{Deng2020831,Pang2023}
      & RL / NN policy \cite{Chaturvedi2024149,Hussien20226442}
      & --
      & Stochastic emulator \cite{Liang20222192,Khalid20241959} \\
    Multi-objective system design
      & NN / RBF \cite{Hirvonen2017323,Wang2017,Schulte2020}
      & $\star$ GP / NN + NSGA \cite{Roth2019214,Han2018622,Aghaei Pour2022303,Reis2025}
      & Retrofit proxy \cite{Thrampoulidis2021}
      & --
      & GP-BO calibration \cite{Chakrabarty2021} \\
    Stochastic / robust planning
      & NN / GP \cite{Lawson201647,Liu20222679,Hu2020}
      & Kriging-NSGA \cite{Lin2023,Pang2023}
      & L2O proxy \cite{Liang20246508}
      & Lagrangian ML \cite{Wu2024391}
      & $\star$ PCE / GP / DQ \cite{Muhlpfordt20194806,Safta20172535,Lu201991827,Chen2023657} \\
    \bottomrule
  \end{tabularx}
\end{table*}
% END inlined table table_T2_task_role_matrix
% BEGIN inlined table table_T5_integration_patterns
\begin{table*}[!t]
  \centering
  \scriptsize
  \caption{Integration patterns for surrogate-assisted energy-system
  optimization. Each row matches one pattern (P1--P5) used in
  Section~\ref{sec:integration}. ``Solver layer'' indicates where the
  surrogate sits in the optimization stack; ``safety hooks'' lists the
  practices that experienced authors use to keep the integration
  robust; ``typical pitfalls'' lists the failure modes that recur in
  the screened pool.}
  \label{tab:T5-integration-patterns}
  \renewcommand{\arraystretch}{1.25}
  \begin{tabularx}{\textwidth}{@{}L{0.04\textwidth} L{0.16\textwidth} L{0.16\textwidth} L{0.18\textwidth} L{0.20\textwidth} L{0.18\textwidth}@{}}
    \toprule
    \# & Pattern name & Surrogate role & Solver layer & Safety hooks & Typical pitfalls \\
    \midrule
    P1 & Replace expensive subroutine
       & Direct functional substitute for an inner optimizer or simulator
       & Inner loop / subproblem
       & Feasibility-restoration layer; constraint-aware embedding; periodic truth-model re-anchoring \cite{Wu2022231,Chen20244723,Chen2022,Liu20242534}
       & Silent constraint violation; OOD operating regimes; surrogate-induced bias of the host objective \\
    P2 & Accelerate search
       & Acquisition / pre-screening of expensive candidates
       & Outer loop / metaheuristic / BO
       & Truth-model check of selected candidates; calibrated posterior for BO; trust-region or infill updates \cite{Nikolaidis2021,Aghaei Pour2022303,Han2018622,Lin2023}
       & Mis-calibrated posterior, uninformative acquisition, wasted truth-budget \\
    P3 & Warm start / solution proxy
       & Predicts (near-)optimal decisions from instance inputs
       & Pre-solve / branch-and-bound seeding
       & Downstream solver as repair step; projection onto feasible set; ensembling for variance reduction \cite{Li2023__2,Chen20244723,Liang20246508,Thrampoulidis2021}
       & Constant-time output is infeasible; warm start dominated by repair cost \\
    P4 & Decomposition support
       & Approximates recourse / pricing / value function
       & Benders / column generation / DP layer
       & Periodic exact cuts; bound tracking on optimality gap; restart on cut staleness \cite{Zhang2025__4,Wu2024391}
       & Surrogate cuts dominate but are not certified; non-convergence of master--subproblem loop \\
    P5 & Uncertainty handling
       & Approximates chance / robust / DRO mapping
       & Reformulation layer
       & Moment or quantile matching; sample-based audit of chance constraints; calibrated uncertainty bands \cite{Muhlpfordt20194806,Safta20172535,Hu2020,Chen2023657}
       & Tail under-coverage; Gaussian-only assumptions; correlation collapse across operating regimes \\
    \bottomrule
  \end{tabularx}
\end{table*}
% END inlined table table_T5_integration_patterns

Five integration patterns (P1--P5) describe \emph{where} a surrogate
sits in the optimization stack, not which surrogate family is used
(that is the subject of Section~\ref{sec:taxonomy}).
In P1 and P3 the surrogate is called whenever the host algorithm needs
a response; in P2 it is called \emph{between} full truth-model runs to
decide which points deserve one; in P4 it supports a coordinating
algorithm that splits one problem into many related subproblems; and in
P5 it translates uncertain inputs into the probabilities or bounds that
enter chance, robust or distributionally robust constraints.
The patterns are not mutually exclusive: a kriging model may screen
candidates in an outer loop (P2) while another surrogate replaces a
network solver inside each accepted candidate (P1), and a
learning-to-optimize proxy (P3) may still invoke a repair step that
resembles P1.
The subsections below follow the same P1--P5 order as
Table~\ref{tab:T5-integration-patterns}; each opens with the defining
contrast to the neighbouring patterns before turning to screened-pool
exemplars.

\subsection{Pattern P1: surrogate replaces an expensive subroutine}
\label{sec:int:replace}

In Pattern P1 the surrogate stands in for a repeated inner evaluation
while the host optimization structure stays unchanged: the mixed-integer
programme, decomposition loop or co-simulation platform still runs, but
one subroutine---power flow, building simulation, device model---is
replaced by a cheap approximation.
The critical distinction from P2 is that P1 substitutes the inner call
itself, not merely the ranking of outer-loop candidates; the
distinction from P3 is that P1 still executes the host solver and only
approximates an embedded function, whereas P3 predicts decisions
directly and may bypass most of the solver altogether.

Li et al.\ train a sparse Gaussian-process surrogate on historical
operating data so a multi-agent reinforcement learner can imitate
coupled power and thermal flow calculations without maintaining
physics-based line and pipe parameters for each microgrid in a
heat--electricity cluster \cite{Li2023__3}. Zheng et al.\ distribute
dispatch of a coupled electricity--heat system with variable
district-heating mass flow: in each alternating iteration the heat
subproblem optimises hydraulic and thermal states through a surrogate
of network heat dynamics before the electricity subproblem is solved on
the reduced model \cite{Zheng20231907}. Gao et al.\ embed a semi-empirical solid-oxide
fuel-cell surrogate inside a MATLAB--TRNSYS co-simulation platform so
multi-objective sizing of a combined cooling, heating and power system
can iterate on component capacities without resolving detailed
electrochemistry at every candidate \cite{Gao2025}. Liu et al.\ embed
piecewise-linear neural surrogates of converter losses---with variable
DC-bus voltage---as linear constraints inside two-stage stochastic
unit commitment for a hybrid AC/DC microgrid \cite{Liu20242534}.

\paragraph{Further cases in Table~\ref{tab:T2-task-role}.}
For economic dispatch, Pang et al.\ embed a general regression neural
network inside an adaptive bat algorithm so large-scale dispatch trials
query a surrogate cost surface instead of the full power-system model
at every iteration \cite{Pang2023}.
Chen et al.\ and Wu et al.\ are also listed under P1 when their
learning-to-optimize models substitute the inner market-clearing or
unit-commitment solve rather than predicting decisions outright
(cf.\ Section~\ref{sec:int:warmstart}) \cite{Chen20244723,Wu2022231}.
Liu et al.\ replace repeated small-signal stability checks inside
security-constrained OPF with an explicit data-driven stability margin
\cite{Liu20223726}.
For capacity decisions, Jahangiri et al.\ apply surrogate models to
map sensitivities of capital costs, demand and carbon pricing on top of
an existing provincial capacity-expansion mixed-integer programme
\cite{Jahangiri20244849}, and
P\'erez-Uresti et al.\ embed rigorous unit surrogates inside a
mixed-integer hydrogen supply-chain expansion model so nonlinear
production and capture units remain tractable across a ten-year plan
\cite{Pérez-Uresti2023}.
At district scale, S\'anchez-Zabala and G\'omez-Acebo and Du et al.\
plug building or network surrogates into district operation studies;
training-data construction is documented in Section~\ref{sec:doe}
\cite{Sánchez-Zabala2024,Du2025}.
Yin et al.\ nest port transport--energy scheduling inside a hybrid
storage capacity-allocation loop where a sampling-based surrogate
accelerates the outer capacity search \cite{Yin2024}.
Roth et al.\ compile hybrid renewable configurations into
mixed-integer sizing programmes through a flexible metamodel layer
\cite{Roth2019214}.
For microgrids, Wang et al.\ benchmark extremal-optimization controller
design against kriging-surrogate alternatives for islanded frequency
control \cite{Wang2017}, and Hussien et al.\ formulate PI tuning
objectives through response-surface metamodels inside a sunflower
optimizer \cite{Hussien20211883}.
In multi-objective community design, Hirvonen et al.\ replace repeated
solar-community simulations with neural metamodelling inside NSGA-II
\cite{Hirvonen2017323}, and Schulte et al.\ substitute reservoir
simulation with Gaussian-process response surfaces when exploring
geothermal well-placement trade-offs \cite{Schulte2020}.
Under stochastic planning, Lawson et al.\ use statistical emulators to
propagate geological and demand uncertainty into transmission planning
decisions \cite{Lawson201647}, Liu et al.\ learn kernel-structured
Gaussian processes for probabilistic load flow \cite{Liu20222679}, and
Hu et al.\ emulate stochastic economic dispatch under wind penetration
with Gaussian-process surrogates \cite{Hu2020}.
Periodic truth-model re-anchoring and constraint-aware embedding, listed
as safety hooks in Table~\ref{tab:T5-integration-patterns}, appear in
the dispatch-proxy studies above \cite{Wu2022231,Chen20244723,Chen2022}.

Across these cases the integration question is whether network,
market or device constraints still hold when the inner evaluation is
approximated (Table~\ref{tab:T5-integration-patterns}, P1 row;
Section~\ref{sec:val:decision}).

\subsection{Pattern P2: surrogate accelerates the search}
\label{sec:int:accelerate}

In Pattern P2 the truth model remains available, but a metamodel
decides which candidates receive a full evaluation.
The surrogate therefore sits in the \emph{outer} loop---inside a genetic
algorithm, Bayesian optimiser or infill strategy---and its job is budget
allocation: predict objective or constraint values cheaply enough to
discard unpromising regions before spending simulation time.
Unlike P1, every accepted solution can still be checked on the
high-fidelity model; unlike P3, P2 does not output a final dispatch or
design vector in one shot.

Reis et al.\ extend NSGA-II for a hybrid solar--wind system with a
symmetry-aware crossover and a surrogate that approximates objective
evaluations in later generations, reducing the number of full
simulation runs needed to reach a diverse Pareto set
\cite{Reis2025}. Salahinezhad et al.\ first run a small set of detailed
trigeneration simulations, fit an RSM to the results, and then use that
surface to compare design choices---solar-collector area, fuel-cell size
and cooling type---without repeating a full hourly plant model at every
trial \cite{Salahinezhad2025}. Granacher and
Mar\'echal screen many alternative plant configurations in
multi-objective process and energy-system design: a classifier estimates
which candidates could still belong to the best cost--performance
trade-off front, and only those are checked with the full nonlinear
model; the active-learning loop that supplies training labels is
documented in Section~\ref{sec:doe:infill}
\cite{Granacher20221573}.

\paragraph{Further cases in Table~\ref{tab:T2-task-role}.}
Nikolaidis and Chatzis use Gaussian-process Bayesian optimization to
propose unit-commitment schedules under volatile renewable behaviour,
checking only the most promising candidates on the full model
\cite{Nikolaidis2021}.
Lin et al.\ and Pang et al.\ pair kriging surrogates with NSGA-II or
bat algorithms for combined economic/emission dispatch on large systems,
using infill rules to refresh the metamodel where it is least accurate
\cite{Lin2023,Pang2023}.
Han et al.\ replace repeated transient-stability simulations in
multi-objective planning of reactive-power compensation with kriging
metamodels inside the evolutionary search \cite{Han2018622}.
Amaral et al.\ adapt metamodel fidelity during simulation-based
capacity-expansion optimisation \cite{Amaral20251156}.
Schulte et al.\ optimise borehole thermal-energy storage layouts with
response-surface metamodels that screen geothermal simulation runs
\cite{Schulte2016104}, and Petrova et al.\ replace repeated aquifer
thermal-interference physics with a random-forest classifier when
screening well placements under geological uncertainty
\cite{Petrova2025}.
For multi-energy and geothermal planning, Han et al., Schulte et al.\
and Aghaei Pour et al.\ combine kriging or Gaussian-process surrogates
with evolutionary or preference-based multi-objective search
\cite{Han2018622,Schulte2020,Aghaei Pour2022303}.
Deng et al.\ accelerate optimal power flow with kriging-assisted
evolutionary algorithms \cite{Deng2020831}.
Roth et al.\ and Reis et al.\ appear again under multi-objective system
design, where metamodels or late-generation surrogates reduce the number
of full hybrid-system evaluations \cite{Roth2019214,Reis2025}.
Truth-model checks of selected candidates and calibrated posteriors for
Bayesian optimisation, listed in Table~\ref{tab:T5-integration-patterns},
motivate the model-management rules in
\cite{Nikolaidis2021,Aghaei Pour2022303,Han2018622,Lin2023}.

The relevant check is search quality per simulation budget
(Table~\ref{tab:T5-integration-patterns}, P2 row), not per-call
feasibility as in P1.

\subsection{Pattern P3: surrogate as warm-start or proxy of the solution map}
\label{sec:int:warmstart}

In Pattern P3 the surrogate maps the inputs that define an operating
case---forecasts, network state, bids or building descriptors---directly
to a decision vector (dispatch, commitment, retrofit measures, or control
trajectory).
The host optimizer, if it is invoked at all, receives a warm start or a
near-feasible candidate rather than calling an expensive inner model
repeatedly: learning-to-optimize and optimisation-proxy architectures
are the dominant forms in the screened pool.
This is the fastest integration mode when it works, but also the most
fragile---small changes in load, network state or market prices can push
the proxy outside its training envelope, which is why repair layers and
feasibility projection (Table~\ref{tab:T5-integration-patterns}, P3 row)
appear so often.

In transmission-level studies, Chen et al.\ train optimisation proxies
for security-constrained economic dispatch that return market-clearing
decisions in milliseconds; their end-to-end architecture adds
differentiable repair layers when the raw network output violates
constraints \cite{Chen2022,Chen20244723}. Wu et al.\ use a convolutional
network to predict which units are on or off, and then compute the
matching hourly dispatch levels in a much smaller follow-up
optimisation step \cite{Wu2022231}. Li et al.\ learn ADMM step updates for
decentralised DC-OPF so agents reach feasible shared solutions faster
\cite{Li2023__2}. (Training-data construction for these proxies is
summarised in Section~\ref{sec:doe}.)

\paragraph{Further cases in Table~\ref{tab:T2-task-role}.}
Thrampoulidis et al.\ map building and climate descriptors to
near-optimal combinations of envelope upgrades and energy-system
choices, replacing repeated building-simulation optimisation runs that
previously took minutes per query; training is covered in
Section~\ref{sec:doe} \cite{Thrampoulidis2021}.
Jalving et al.\ train neural market surrogates that predict clearing
outcomes from integrated energy system bids and can seed or warm-start
market-aware design iterations \cite{Jalving2023}.
Wang et al.\ fit a polynomial response surface to integrated energy
system operating costs so the scheduler queries a cheap proxy objective
\cite{Wang2020202933}.
Reynolds et al.\ train an artificial neural network on EnergyPlus runs
to supply fitness values to a multi-objective genetic set-point scheduler
for district-heating-compatible building control \cite{Reynolds2017704}.
Chaturvedi et al.\ embed a neural-network surrogate in a DC microgrid
control loop: it predicts power losses and uneven load sharing from the
local voltage-sharing settings, so reinforcement learning can adjust
those settings online without repeated full system simulation
\cite{Chaturvedi2024149}. Hussien et al.\ tune PI controllers for autonomous microgrid operation
with a coot-bird metaheuristic that treats each closed-loop trial as
an expensive evaluation \cite{Hussien20226442}.
Liang et al.\ accelerate joint chance-constrained unit commitment with
learning-to-optimize warm starts for the robust inner programme
\cite{Liang20246508}.
Downstream repair layers and projection onto feasible sets
(Table~\ref{tab:T5-integration-patterns}) are explicit in
\cite{Chen20244723,Liang20246508,Thrampoulidis2021}.

Multi-carrier evidence for one-shot proxies is still thin: most
coupled-carrier papers in the pool substitute inner evaluations (P1)
or metamodels within search (P2) rather than predicting full joint
schedules directly.

\subsection{Pattern P4: surrogate enables decomposition}
\label{sec:int:decomp}

In Pattern P4 the optimizer does not solve one monolithic model at once.
It splits the problem into a coordinating step and one or more smaller
subproblems, then iterates: the coordinator proposes, each subproblem
responds, and the coordinator updates its plan.
The surrogate approximates what those subproblems would answer---for
example, how a virtual power plant would schedule given a market price,
which network constraints are likely to matter, or which penalty
weights make the subproblems easy to solve---so the coordinator can move
forward without running every subproblem to full accuracy at every
iteration.
P4 therefore differs from P1 because the cheap model serves this
back-and-forth coordination, not a single repeated physics or market
call inside one solver, and from P3 because the output still feeds an
ongoing decomposition loop rather than the final operating decision
directly.

Coupled-carrier examples are rare in the screened pool. Cao et al.\
study a market operator and several virtual power plants that trade
electricity, heat and carbon: the operator proposes prices, each plant
chooses its schedule in response, and a Kriging model learns how each
plant would react to a given price so many price options can be compared
without rerunning the full plant optimisation every time
\cite{Cao2023__2}. Zhang et al.\ speed up hour-by-hour simulation of
generation, the transmission grid and storage operating together by
splitting the model into discrete operating choices and continuous
hourly decisions; a Gaussian-process model learns from past runs which
constraints are unlikely to matter and leaves them out of the next solve
\cite{Zhang2025__4}. Wu et al.\ decompose hourly unit commitment into
smaller hourly subproblems and train deep networks to predict penalty
weights that make those subproblems easier to solve, so the overall
on/off schedule can be found faster without learning the full
combinatorial problem at once \cite{Wu2024391}. All
three report lower runtime for a fixed solution target; certifying how
much accuracy is lost by the learned approximations remains open
(Section~\ref{sec:challenges}).
Periodic exact cuts and bound tracking on the optimality gap
(Table~\ref{tab:T5-integration-patterns}) are the integration safeguards
emphasised in \cite{Zhang2025__4,Wu2024391}.

\subsection{Pattern P5: surrogate handles uncertainty}
\label{sec:int:uncertainty}

In Pattern P5 the surrogate estimates how uncertain inputs, e.g.\ wind
and load forecasts and market prices, translate into uncertain outcomes
such as operating costs. The optimiser uses those probabilities when it
must stay within safety or reliability limits under uncertainty. P5
surrogates provide those estimates without running a large random sample
of full model evaluations at every optimisation step.
Unlike P1--P4, the main question is not ``how cheap is this candidate?''
but ``how likely is a limit to be violated?''; validation therefore
focuses on rare bad cases and constraint satisfaction under hold-out
scenarios rather than mean prediction error alone.

Dong et al.\ model renewable-generation and load uncertainty for a
renewable-energy-source combined cooling, heating and power system with
Gaussian-process regressions that return predictive distributions, then
embed those distributions in a robust model-predictive rolling
optimisation layer \cite{Dong2022}. Chaudhry et al.\ quantify market
and demand risk for fifth-generation district heating and cooling
networks through polynomial-chaos Kriging expansions that substitute
repeated high-fidelity network evaluations inside robust design and
operational optimisation \cite{Chaudhry2024}. In electricity-only
studies, M\"uhlpfordt et al.\ use polynomial-chaos expansions for joint
chance-constrained AC-OPF \cite{Muhlpfordt20194806}, Safta et al.\
replace Monte Carlo loops in stochastic economic dispatch
\cite{Safta20172535}, and Chen et al.\ learn quantile surfaces for
chance-constrained distribution-network OPF when impedance data are
partially unknown \cite{Chen2023657}.

\paragraph{Further cases in Table~\ref{tab:T2-task-role}.}
Hu et al.\ construct Gaussian-process emulators of stochastic economic
dispatch after Karhunen--Lo\`eve dimension reduction, replacing Monte
Carlo loops at scale \cite{Hu2020,Hu2021671}.
Dong et al.\ reformulate chance-constrained day-ahead dispatch of an
integrated energy system from sparse polynomial-chaos expansions of
power-flow responses; sampling for the expansions is treated in
Section~\ref{sec:doe} \cite{Dong2023}.
Wang et al.\ link photovoltaic, storage and electrolysis design variables
to techno-economic objectives through a machine-learning surrogate inside
multi-objective hydrogen-system optimisation \cite{Wang2024}.
Fu et al.\ learn stochastic distribution-network expansion responses
with statistical machine-learning models that support dimension-reduced
planning under dynamic correlation \cite{Fu20202904}.
Kim et al.\ propagate wind-resource uncertainty into probabilistic
expansion screening \cite{Kim2021}.
Simonson et al.\ train physics-informed Gaussian-process models of
district-heating pipe dynamics for stochastic nonlinear predictive
control \cite{Simonson202036}.
Liang et al.\ use compressive sensing to solve stochastic energy-router
storage management without repeated scenario expansions
\cite{Liang20222192}.
Khalid et al.\ schedule hybrid energy networks under renewable
uncertainty with radial-basis-function emulators inside a metaheuristic
search \cite{Khalid20241959}.
Chakrabarty et al.\ calibrate coupled building--HVAC parameters through
scalable Bayesian optimisation when each calibration trial would
otherwise require an expensive coupled simulation solve
\cite{Chakrabarty2021}.
Lu et al.\ nest sparse-grid stochastic collocation inside economic
dispatch under wind and photovoltaic uncertainty \cite{Lu201991827}.
Koirala et al.\ propagate photovoltaic uncertainty into hosting-capacity
chance constraints through general polynomial chaos, avoiding repeated
power-flow solves \cite{Koirala20242284}.
Schulte et al.\ explore geothermal design trade-offs through Gaussian-process
response surfaces over an ensemble of geological models
\cite{Schulte2020}.
Moment matching and sample-based audits of chance constraints
(Table~\ref{tab:T5-integration-patterns}) motivate the validation focus
in \cite{Muhlpfordt20194806,Safta20172535,Hu2020,Chen2023657}.

Integration-specific validation here targets tail behaviour and
constraint satisfaction under hold-out uncertainty scenarios
(Table~\ref{tab:T5-integration-patterns}, P5 row;
Section~\ref{sec:val:feasibility}), not mean error alone.
